\chapter{Estimation}\label{ch:iv:estimation}

``How many messages does the US options market's consolidated feed publish in its busiest ten milliseconds of
a month?'' The candidate who answers ``between two hundred thousand and two million, most likely around six
hundred thousand'' and shows the chain (a few thousand series on a busy underlying, seventeen exchanges
quoting each, several moves of the underlying in the window) has answered better than one who happens to
know the published figure and cannot say how it comes about. An estimation question measures two things: the
ability to decompose an unknown quantity into knowable pieces, and the honesty to say how wrong the result
might be.

\section{The Fermi chain}

\begin{definition}[Fermi estimate]\label{def:iv:estimation:fermi}
A \emph{Fermi estimate}\index{Fermi estimate} is an order-of-magnitude estimate of an unknown quantity
obtained by writing it as a product (or sum) of factors each of which can be estimated from common knowledge,
and multiplying the estimates.
\end{definition}

\begin{method}[Building a Fermi chain]\label{met:iv:estimation:chain}
\begin{enumerate}
\item Write the quantity as a product of three to six factors, each with a unit, so that the units cancel to
the answer's unit.
\item Estimate each factor from something you know: a population, a price, a time, a size.
\item Where only bounds come to mind, use their geometric mean: between 200 and 20\,000 means about 2\,000.
\item Multiply in powers of ten first, then the mantissas.
\item Give an interval (next section) and check against a second chain built from different factors.
\end{enumerate}
\end{method}

\begin{example}[The hook's chain]\label{ex:iv:estimation:opra}
Suppose (these are assumptions, stated as such) that a heavily traded underlying has about 5\,000 listed option
series, that 17 exchanges quote each series, and that in a burst the underlying moves about ten times in ten
milliseconds across the names that move together, each move requoting every series once: $5\,000 \times 17
\times 10 = 850\,000$ messages. The published peak for July 2026 was 0.89 million messages in ten
milliseconds (\cref{dat:iv:estimation:anchors}). The chain's closeness is partly luck; its structure, a
product of breadth, venues and event rate, is what the interviewer wants.
\end{example}

Errors in a chain multiply, so they add on the log scale, and independent errors partly cancel.

\begin{proposition}[How uncertainty grows along a chain]\label{prop:iv:estimation:logsd}
If each of $k$ factors is estimated with an independent log error of standard deviation $s_i$, the log error of
the product has standard deviation $\big(\sum s_i^2\big)^{1/2}$. With four factors each at $s = 0.35$ (each
within a factor of about 1.8 with 90\% probability), the product has $s = 0.70$, and a 90\% interval runs from a
factor $e^{1.645 \times 0.7} \approx 3.2$ below to 3.2 above the estimate.
\end{proposition}

\begin{proof}
The log of a product is the sum of the logs; variances of independent errors add. The interval factor is
$\exp(z_{0.95}\, s)$ for a lognormal error.
\end{proof}

\section{Intervals}

\begin{definition}[Calibrated interval]\label{def:iv:estimation:calibrated}
A stated $c$ interval for an unknown quantity is a \emph{calibrated interval}\index{calibrated interval} if,
over many such statements, the true value falls inside a fraction $c$ of them.
\end{definition}

Calibration is a property of the estimator, not of one interval, and people are usually overconfident: in the
studies surveyed by Lichtenstein, Fischhoff and Phillips (1982), intervals stated with high confidence
contained the truth much less often than stated. \Cref{fig:iv:estimation:calib} shows the arithmetic: an
estimator whose log errors have standard deviation 1 but who builds intervals as if it were 0.5 hits a stated
90\% interval only 59\% of the time.

\begin{omfigure}
\begin{tikzpicture}
\begin{axis}[width=9.5cm,height=5.4cm,xlabel={\small stated coverage (\%)},ylabel={\small actual coverage (\%)},
  xmin=50,xmax=95,ymin=20,ymax=100,
  tick label style={font=\footnotesize},grid=major,grid style={gray!20},table/col sep=comma,
  legend cell align=left,legend style={font=\scriptsize,at={(0.5,-0.32)},anchor=north,
  /tikz/every even column/.append style={column sep=8pt}},legend columns=3]
\addplot[omDef,thick] table[x=stated,y=calibrated]{figdata/interviews/09-estimation/calib.csv};
\addplot[omThm,thick,dashed] table[x=stated,y=overconfident]{figdata/interviews/09-estimation/calib.csv};
\addplot[only marks,mark=o,mark size=1.6pt,omThm] table[x=stated,y=sim_over]{figdata/interviews/09-estimation/calib.csv};
\legend{calibrated,overconfident (exact),overconfident (simulated)}
\end{axis}
\end{tikzpicture}

\omcaption{Actual against stated coverage of central intervals when log errors have standard deviation 1: a
calibrated estimator (intervals built with 1) lies on the diagonal; an overconfident one (intervals built
with 0.5) falls far below it. Circles: 20\,000 simulated estimates. Data:
\texttt{fig\_iv\_calib.py}.}\label{fig:iv:estimation:calib}
\end{omfigure}

An interval is scored by a proper scoring rule that rewards narrowness and penalises misses in proportion to
their size. The interval score of Gneiting and Raftery (2007) for a central $(1-\alpha)$ interval $[l, u]$ and
outcome $x$ is $(u - l) + \tfrac2\alpha(l - x)\mathbf 1_{x<l} + \tfrac2\alpha(x - u)\mathbf 1_{x>u}$: an
honest interval minimises its expectation. The Brier score of One Quant Book~16, chapter~26, plays the same role
for probabilities of events.

\section{Market estimation questions}

Market questions are Fermi questions with a trading context: volumes, notionals, messages, bytes, capacities.
The factors come from the reader's knowledge of the markets of One Quant Books 1 to~3 and the systems of Books
13 and~14, and the published anchors for a few of them are in \cref{dat:iv:estimation:anchors}. Two habits
distinguish good answers. The first is to separate the count of things from the activity per thing (series
times quote rate; accounts times trades per account). The second is to state the peak-to-average ratio
explicitly when a capacity is asked, because systems are sized for the burst, not the day.

\begin{dated}{2026-09}{Anchors for market estimates}\label{dat:iv:estimation:anchors}
Global over-the-counter FX turnover was USD 9.6 trillion a day in April 2025 (BIS Triennial Survey). The US
options consolidated feed (OPRA) published peaks in July 2026 of 63.9 million messages in a second and 0.89
million in ten milliseconds, and its capacity projection for July 2026 was 4.403 gigabits per 100 ms on one
stream. CME Group's exchanges averaged 28.1 million contracts a day in 2025. The world's population was 8.2
billion in 2024 (UN World Population Prospects 2024).
\end{dated}

\section{Checking an estimate with a second chain}

A second chain that shares no factor with the first is the best check an estimate can have. If the two agree
within their intervals, combine them by a weighted geometric mean (weights inversely proportional to the log
variances); if they disagree by more, one of the chains has a wrong factor, and finding it is worth more than
averaging (\cref{iq:iv:estimation:11}).

\section{Worked answers}

\begin{example}[Multiplying in logarithms]\label{ex:iv:estimation:logs}
A chain ends in $3.7 \times 10^3 \times 2.2 \times 10^4 \times 0.45$. In base-10 logarithms, $\log 3.7 \approx 0.57$,
$\log 2.2 \approx 0.34$ and $\log 0.45 \approx -0.35$, so the product has logarithm $0.57 + 0.34 - 0.35 + 7 = 7.56$, and
$10^{0.56} \approx 3.6$: about $3.6 \times 10^7$. The exact product is $3.663 \times 10^7$. Estimators who work in
logarithms carry the exponent separately, never lose a zero, and can state the uncertainty of each factor in the same
units.
\end{example}

\begin{example}[A chain, its interval and a second chain]\label{ex:iv:estimation:orders}
\emph{``How many order messages (new orders, modifications, cancels) does an options market maker send on a
trading day? Give a 90\% interval.''}

\emph{First chain.} The firm quotes, say, 2\,000 option series actively on one venue; each quote is updated about 5 times a
second on average, as the underlying moves; the session has 6.5 hours, $23\,400$ seconds. The product is $2\,000 \times 5
\times 23\,400 \approx 2.3 \times 10^8$. The instruments and the update rate are each uncertain to a factor of 3 at 90\%;
in natural logarithms each has a standard deviation of $\ln 3/1.645 \approx 0.67$, together $0.67\sqrt2 \approx 0.94$,
and the 90\% factor is $e^{1.645 \times 0.94} \approx 4.7$. Interval: about $5 \times 10^7$ to $1.1 \times 10^9$.

\emph{Second chain, from the pipe.} If the firm's order-entry traffic on that venue averages 10 megabits a second and a
message is about 50 bytes, it sends $10^7/(8 \times 50) = 25\,000$ messages a second, $5.85 \times 10^8$ a day. The two
chains share no factor and agree within a factor of 2.5, inside the interval; if they had disagreed by a factor of 20,
the right answer would have been to find out which assumption was wrong, not to average them (the
previous section).
\end{example}

\section{Question bank}

\begin{interviewq}[$\star$ \iqroles{trader} \iqfirm{market maker}]\label{iq:iv:estimation:1}
How many litres of coffee does a trading floor of 400 people drink in a year? Show the chain.
\end{interviewq}

\begin{interviewq}[$\star$ \iqroles{trader, developer} \iqfirm{any}]\label{iq:iv:estimation:2}
How many seconds are there in a year, and how many seconds of regular trading does a US equity market have in
a year?
\end{interviewq}

\begin{interviewq}[$\star$ \iqroles{trader, researcher} \iqfirm{any}]\label{iq:iv:estimation:3}
You are sure a quantity lies between 200 and 20\,000 and have no other information. What single number do you
give, and why not 10\,100?
\end{interviewq}

\begin{interviewq}[$\star$ \iqroles{trader, bank} \iqfirm{bank}]\label{iq:iv:estimation:4}
Give a 90\% interval for the daily turnover of the global FX market, with a chain.
\end{interviewq}

\begin{interviewq}[$\star\star$ \iqroles{developer, trader} \iqfirm{market maker}]\label{iq:iv:estimation:5}
Give a 90\% interval for the largest number of messages the US options consolidated feed publishes in any ten
milliseconds of a month. What would you ask before building a system to receive it?
\end{interviewq}

\begin{interviewq}[$\star\star$ \iqroles{trader} \iqfirm{proprietary firm}]\label{iq:iv:estimation:6}
Estimate the number of futures and options contracts traded on an average day on the exchanges of the largest
US derivatives exchange group. Give an interval.
\end{interviewq}

\begin{interviewq}[$\star\star$ \iqroles{trader, researcher} \iqfirm{any}]\label{iq:iv:estimation:7}
How many people alive today have their birthday today?
\end{interviewq}

\begin{interviewq}[$\star\star$ \iqroles{developer} \iqfirm{proprietary firm}]\label{iq:iv:estimation:8}
Estimate the storage needed for one day of top-of-book updates for all US-listed stocks, uncompressed.
\end{interviewq}

\begin{interviewq}[$\star\star$ \iqroles{researcher, risk} \iqfirm{systematic fund}]\label{iq:iv:estimation:9}
Over the last month you gave ten 90\% intervals and five contained the true value. How surprising is that if
you were calibrated? What do you change?
\end{interviewq}

\begin{interviewq}[$\star\star\star$ \iqroles{researcher, trader} \iqfirm{any}]\label{iq:iv:estimation:10}
A chain has four independent factors, each of which you believe is within a factor of 2 of the truth with 90\%
probability. Within what factor is the product, at 90\%? Why is it less than $2^4$?
\end{interviewq}

\begin{interviewq}[$\star\star\star$ \iqroles{researcher} \iqfirm{systematic fund}]\label{iq:iv:estimation:11}
Two independent chains give $10^6$ (log standard deviation 0.7) and $10^7$ (log standard deviation 1.0). What
single estimate do you report, with what uncertainty? When would you refuse to combine them?
\end{interviewq}

\begin{interviewq}[$\star\star\star$ \iqroles{researcher, risk} \iqfirm{bank}]\label{iq:iv:estimation:12}
Two candidates give 90\% intervals for daily FX turnover: $[5, 12]$ trillion and $[2, 5]$ trillion dollars. The
answer is 9.6 trillion. Score both with the interval score, and explain why the rule makes honest intervals
the best strategy.
\end{interviewq}

\begin{interviewq}[$\star\star\star$ \iqroles{developer} \iqfirm{market maker}]\label{iq:iv:estimation:13}
At a peak of 63.9 million messages a second and an average of 40 bytes a message, what bandwidth does the
options feed need, in gigabits a second? Compare with a published capacity of 4.403 gigabits per 100
milliseconds, and say what your estimate leaves out.
\end{interviewq}

\begin{omsources}
\item S.~Lichtenstein, B.~Fischhoff and L.~D.~Phillips, ``Calibration of probabilities: the state of the art to
1980'', in D.~Kahneman, P.~Slovic and A.~Tversky (eds), \emph{Judgment under Uncertainty}, Cambridge University
Press, 1982, 306--334.
\item T.~Gneiting and A.~E.~Raftery, ``Strictly proper scoring rules, prediction, and estimation'',
\emph{Journal of the American Statistical Association} 102(477), 2007, 359--378.
\item BIS, Triennial Central Bank Survey 2025; OPRA, key operating metrics (August 2026) and capacity
projections (September 2025); CME Group, Form 10-K for 2025; UN, \emph{World Population Prospects 2024}.
\end{omsources}
